\subsubsection{Relação entre requisitos, funções e restrições}
\label{sec:requisitos-funcoes}

A revisão de 15 de outubro de 2026 distingue o que o produto deve fazer, o desempenho exigido e as condições que restringem a solução. Requisitos funcionais e de desempenho podem aparecer na mesma especificação: amortecer impacto descreve um efeito; limitar o pico de aceleração estabelece seu desempenho. Condições ambientais, durabilidade, interfaces e custo também integram o conjunto técnico \citep{nasarequisitos2023,nasadecomposicao2023}.

A análise parte do efeito necessário para cumprir cada requisito. Quando esse efeito já existe, o requisito qualifica a função; quando uma ação adicional é necessária, deriva-se uma função de suporte. Quando se trata de um limite econômico ou geométrico, registra-se a restrição e sua aplicação. A especificação-meta permanece verificável em todos os casos. Expressar uma função por verbo e objeto não basta: deve existir ação, efeito ou capacidade necessária, com fronteira e vínculo identificáveis. Funções também abrangem suporte e comunicação, além de transformações de energia ou material.

A aba Função Global conserva 23 vínculos, V01--V23, associados aos 23 RP, e mantém um catálogo separado com 17 funções. As classes técnica, estrutural, ergonômica, interativa, semântica, sintática e adicional caracterizam as funções; não indicam se o requisito é funcional ou não funcional. Uma função pode responder a vários requisitos, e um requisito pode depender de várias funções. Os IDs numéricos preservados no catálogo não formam uma sequência completa, pois identificadores redundantes ou usados para atributos foram descontinuados.

\begin{table}[H]
\caption{Encaminhamento dos requisitos de qualidade, suporte e restrição}
\label{tab:requisito-funcao}
\centering\scriptsize
\begin{tabularx}{\textwidth}{L{2.4cm}YL{3.3cm}}
\toprule\textbf{Requisito}&\textbf{Vínculo ou desdobramento funcional}&\textbf{Critério que permanece}\\\midrule
RP12: uniformidade&Qualifica F01, Amortecer impacto; não cria outra função de amortecimento&Coeficiente de variação do pico $\leq10\%$\\
RP13: umidade&Qualifica F11, Suportar carga, após condicionamento; bloquear água só será função auxiliar se necessário ao princípio&Perda de capacidade $\leq12{,}5\%$ após 48\,h em câmara úmida\\
RP14: volume&Restringe F03/F08 e a geometria carregada. F14, Recolher partes, atua no retorno vazio&Medir carregado e vazio; recolhimento não comprova o limite carregado\\
RP15: manutenção&F15, Prover acesso; F17, Liberar partes; F22, Unir partes, viabilizam a operação humana&Ciclo de limpeza, inspeção e reposição $\leq5$\,min\\
RP18: vida&Qualifica proteção e suporte ao reuso; não é nova transformação nem reaproveitamento entre gerações&$\geq20$ ciclos até falha funcional\\
RP19: reciclagem&F17/F19 viabilizam liberação e separação material; composição e cadeia regional permanecem critérios&$\geq95\%$ da massa com cadeia estabelecida\\
RP20: modularidade&Restringe arquitetura e interfaces; acesso e uniões precisam permitir troca independente&Demonstrar $\geq3$ módulos substituíveis\\
RP23: custo&Restringe todas as funções e soluções; não gera bloco físico autônomo&Teto provisório de R\$\,15/ciclo e comparação econômica completa\\
\bottomrule\end{tabularx}
\fonte{Aplicação da análise funcional aos RP e às metas de trabalho da equipe. Valores ainda sujeitos à validação.}
\end{table}

\begin{table}[H]
\caption{Catálogo de funções e capacidades de suporte}
\label{tab:catalogo-funcoes}
\centering\scriptsize
\begin{tabularx}{\textwidth}{L{0.8cm}L{3.5cm}Y}
\toprule\textbf{ID}&\textbf{Função}&\textbf{Efeito ou capacidade requerida}\\\midrule
F01&Amortecer impacto&Atenuar choque e vibração antes da energia chegar à cerâmica\\
F02&Imobilizar louça&Conservar o posicionamento da carga na fronteira definida\\
F03&Portar louça&Transferir a carga entre posições com interfaces de movimentação\\
F04&Expor condição&Tornar a região crítica da cerâmica acessível à inspeção\\
F05&Prover pega&Disponibilizar interface para a mão enluvada aplicar força\\
F06&Reter fragmentos&Conter fragmentos também durante abertura e retirada\\
F07&Acomodar variação&Adaptar o posicionamento à faixa dimensional sem trocar partes no ajuste\\
F08&Compatibilizar carga&Compatibilizar geometria e interfaces com palete e transporte\\
F09&Estabilizar conjunto&Manter o conjunto carregado apoiado contra tombamento ou movimento indesejado\\
F11&Suportar carga&Transmitir compressão pelo caminho resistente do conjunto\\
F14&Recolher partes&Mudar a configuração vazia para reduzir o volume de retorno\\
F15&Prover acesso&Viabilizar acesso às regiões de limpeza, inspeção e reposição\\
F16&Comunicar uso&Comunicar ação, orientação, fechamento e advertências\\
F17&Liberar partes&Desfazer uniões para abertura, manutenção e fim de vida\\
F19&Separar materiais&Viabilizar frações separáveis para destinação no fim de vida\\
F21&Exibir identidade&Exibir identificador permanente e único da unidade\\
F22&Unir partes&Estabelecer as uniões que mantêm o conjunto montado e contido\\
\bottomrule\end{tabularx}
\fonte{Catálogo da aba Função Global. As funções de suporte não pressupõem componentes ou módulos exclusivos.}
\end{table}

F10 foi reunida em F16; F12 foi vinculada a F01; RP13 passou a qualificar F11. Os antigos F18 e F20 foram substituídos pelos vínculos de durabilidade e arquitetura às funções necessárias. O rótulo F23 deixou de ser ID funcional; RP23 permanece restrição de custo. Na morfológica, Prover acesso explicita o suporte antes denominado Permitir manutenção. Recolher partes explicita a ação no retorno antes denominada Limitar volume. A comunicação de uso acompanha a identidade, pois um serial ou RFID sozinho não orienta fechamento e manuseio.

FAST e estruturas alternativas usam os IDs do catálogo e mantêm atributos nas condições de verificação. A morfológica acrescenta os IDs associados aos princípios existentes. A MIM continua avaliando os mesmos 13 portadores e os mesmos 156 julgamentos; foram atualizados os vínculos funcionais, sem criar peças ou módulos para uniformidade, vida ou custo. Nenhuma dessas correções representa resultado de ensaio ou fechamento da seleção.
